\documentclass[11pt]{article}

\usepackage[margin=1in]{geometry}
\usepackage{booktabs}
\usepackage{enumitem}
\usepackage{hyperref}
\usepackage{xcolor}
\hypersetup{colorlinks=true, linkcolor=blue, urlcolor=blue, citecolor=blue}

\title{Supplementary Material: Dataset Documentation\\[0.5em]
\large CEI: A Benchmark for Evaluating Pragmatic Reasoning in Language Models}

\author{Chun et al.\\Kenyon College, Gambier, OH, USA}
\date{February 2026}

\begin{document}
\maketitle

\noindent This document provides the required DMLR dataset documentation for the CEI Benchmark as a standalone supplementary PDF, following the Datasheets for Datasets framework \cite{gebru2021datasheets}.

\section{Motivation}

\begin{itemize}[leftmargin=*]
    \item \textbf{Purpose:} Evaluate emotion inference from pragmatically complex (indirect) speech in language models. CEI tests whether models can move beyond surface semantics to infer what a speaker actually feels.
    \item \textbf{Creators:} Academic research team at Kenyon College, led by Jon Chun.
    \item \textbf{Funding:} None; completed as part of university coursework (IPHS 391, Fall 2025).
    \item \textbf{Prior uses:} This is the first public release. Zero-shot baselines across 7 LLMs are reported in the accompanying paper.
\end{itemize}

\section{Composition}

\begin{itemize}[leftmargin=*]
    \item \textbf{Instances:} 300 scenarios, each consisting of a situational context (2--4 sentences), labeled speaker and listener roles with power relation, and a pragmatically ambiguous utterance.
    \item \textbf{Labels per scenario:} 3 independent annotations, each providing:
    \begin{itemize}
        \item Primary emotion from Plutchik's 8 basic emotions (joy, trust, fear, surprise, sadness, disgust, anger, anticipation)
        \item Valence-Arousal-Dominance (VAD) ratings on 7-point scales mapped to $[-1.0, +1.0]$
        \item Annotator confidence rating
    \end{itemize}
    \item \textbf{Gold labels:} Derived via majority vote; three-way splits adjudicated by expert meta-annotator.
    \item \textbf{Pragmatic subtypes:} 5 (sarcasm/irony, mixed signals, strategic politeness, passive aggression, deflection/misdirection), 60 scenarios each.
    \item \textbf{Power relations:} 3 (peer-to-peer: 217, high-to-low: 61, low-to-high: 22).
    \item \textbf{Social settings:} Workplace, family, social/friendship, service encounters.
    \item \textbf{Splits:} Predefined stratified train/val/test (70/15/15) balanced by subtype and power relation, generated with seed=42. Available in \texttt{reports/cei2026/splits.json}.
    \item \textbf{Sensitive content:} None. All scenarios are synthetic (researcher-authored). No real conversations, personal data, or copyrighted materials are included.
    \item \textbf{Personally identifiable information:} None.
\end{itemize}

\section{Collection Process}

\begin{itemize}[leftmargin=*]
    \item \textbf{Scenario generation:} Template-based synthesis with expert curation. 600 initial scenarios were narrowed to 300 after manual review for quality, clarity, and subtype balance.
    \item \textbf{Annotation:} 15 undergraduate annotators from an interdisciplinary course (IPHS 391) at Kenyon College, all fluent English speakers. Each annotator was assigned all 60 scenarios of a single subtype (3 annotators per subtype).
    \item \textbf{Training:} Two 20-minute in-class explanations plus written guidelines (1-page cheatsheet and 3-page detailed guideline document).
    \item \textbf{Annotation tool:} Label Studio (open-source).
    \item \textbf{Timeline:} November--December 2025. Each annotator completed their 60 scenarios in a single session of approximately 50--70 minutes.
    \item \textbf{Quality control:} 4-level pipeline: (1) schema validation, (2) statistical consistency checks (straight-lining, timing outliers, self-contradiction), (3) agreement analysis (Fleiss' $\kappa$ with bootstrap CIs), (4) expert adjudication of flagged items (15.7\% of dataset).
    \item \textbf{Ethics:} IRB-exempt (Category 2). Annotators received course credit and provided informed consent for publication of their annotation labels.
    \item \textbf{Consent:} All annotators were informed that their labels (but not free-text explanations) would be released publicly. Co-authorship terms were communicated before annotation began.
\end{itemize}

\section{Preprocessing}

\begin{itemize}[leftmargin=*]
    \item Annotations exported from Label Studio as JSON, then processed through the 4-level QA pipeline.
    \item 2.6\% of annotations flagged at Level 2 (statistical consistency); none met exclusion thresholds.
    \item 15.7\% of scenarios (47/300) required Level 4 expert adjudication.
    \item Ground truth established via majority vote + expert override for three-way splits.
    \item VAD ratings mapped from 7-point integer scales to $[-1.0, +1.0]$ continuous range.
\end{itemize}

\section{Uses}

\begin{itemize}[leftmargin=*]
    \item \textbf{Intended uses:} Evaluating LLM pragmatic reasoning capabilities; benchmarking contextual emotion inference; studying inter-annotator agreement on subjective tasks; training emotion inference models (with appropriate caveats about dataset size).
    \item \textbf{Not recommended:} High-stakes decision-making without human oversight; surveillance of employee sentiment; deployment as sole ground truth for emotion classification in clinical or legal settings.
\end{itemize}

\section{Author Responsibility Statement}

The authors bear all responsibility in case of violation of rights, and confirm that:
\begin{itemize}[leftmargin=*]
    \item The dataset is released under the CC-BY-4.0 license.
    \item All accompanying code is released under the MIT license.
    \item All scenarios are synthetic and do not contain personal data, copyrighted materials, or content from real conversations.
    \item The collection process was reviewed and deemed IRB-exempt (Category 2).
\end{itemize}

\section{Distribution and Hosting}

\begin{itemize}[leftmargin=*]
    \item \textbf{License:} CC-BY-4.0 (data), MIT (code)
    \item \textbf{Primary access:} \url{https://github.com/jon-chun/cei-tom-dataset-public}
    \item \textbf{Persistent archive:} Zenodo (DOI: \texttt{10.5281/zenodo.18528706})
    \item \textbf{Machine-readable metadata:} Croissant format (auto-generated by HuggingFace Hub)
    \item \textbf{File formats:} CSV (annotations), JSON (Label Studio exports, splits), YAML (configuration)
\end{itemize}

\section{Maintenance Plan}

\begin{itemize}[leftmargin=*]
    \item \textbf{Versioning:} Semantic versioning (current: v1.0.0). All corrections logged in a changelog.
    \item \textbf{Support commitment:} Minimum 5 years post-publication.
    \item \textbf{Hosting:} GitHub (primary, actively maintained) + Zenodo (persistent archive with DOI).
    \item \textbf{Contact:} \texttt{chunj@kenyon.edu}
    \item \textbf{Issue tracking:} GitHub Issues on the repository.
\end{itemize}

\begin{thebibliography}{1}
\bibitem{gebru2021datasheets}
T.~Gebru, J.~Morgenstern, B.~Vecchione, J.~W. Vaughan, H.~Wallach, H.~{Daum\'{e} III}, and K.~Crawford.
\newblock Datasheets for datasets.
\newblock {\em Communications of the ACM}, 64(12):86--92, 2021.
\end{thebibliography}

\end{document}
